\documentclass[10pt,a4paper]{amsart}
\usepackage{amsmath,amssymb,amsthm,graphicx}
\usepackage{verbatim}
\newenvironment{reference}{\comment}{\endcomment}
%\newenvironment{reference}{}{}
\newenvironment{slogan}{\comment}{\endcomment}
\newenvironment{history}{\comment}{\endcomment}

% For commutative diagrams we use Xy-pic
\usepackage[all]{xy}

% We use 2cell for 2-commutative diagrams.
\xyoption{2cell}
% Theorem environments.
%
\theoremstyle{plain}
\newtheorem{theorem}[subsection]{Theorem}
\newtheorem{proposition}[subsection]{Proposition}
\newtheorem{lemma}[subsection]{Lemma}

\theoremstyle{definition}
\newtheorem{definition}[subsection]{Definition}
\newtheorem{example}[subsection]{Example}
\newtheorem{exercise}[subsection]{Exercise}
\newtheorem{situation}[subsection]{Situation}

\theoremstyle{remark}
\newtheorem{remark}[subsection]{Remark}
\newtheorem{remarks}[subsection]{Remarks}

\numberwithin{equation}{subsection}

% Macros
%
\def\lim{\mathop{\mathrm{lim}}\nolimits}
\def\colim{\mathop{\mathrm{colim}}\nolimits}
\def\Spec{\mathop{\mathrm{Spec}}}
\def\Hom{\mathop{\mathrm{Hom}}\nolimits}
\def\Ext{\mathop{\mathrm{Ext}}\nolimits}
\def\SheafHom{\mathop{\mathcal{H}\!\mathit{om}}\nolimits}
\def\SheafExt{\mathop{\mathcal{E}\!\mathit{xt}}\nolimits}
\def\Sch{\mathit{Sch}}
\def\Mor{\mathop{\mathrm{Mor}}\nolimits}
\def\Ob{\mathop{\mathrm{Ob}}\nolimits}
\def\Sh{\mathop{\mathit{Sh}}\nolimits}
\def\NL{\mathop{N\!L}\nolimits}
\def\CH{\mathop{\mathrm{CH}}\nolimits}
\def\proetale{{pro\text{-}\acute{e}tale}}
\def\etale{{\acute{e}tale}}
\def\QCoh{\mathit{QCoh}}
\def\Ker{\mathop{\mathrm{Ker}}}
\def\Im{\mathop{\mathrm{Im}}}
\def\Coker{\mathop{\mathrm{Coker}}}
\def\Coim{\mathop{\mathrm{Coim}}}

% Boxtimes
%
\DeclareMathSymbol{\boxtimes}{\mathbin}{AMSa}{"02}

%
% Macros for moduli stacks/spaces
%
\def\QCohstack{\mathcal{QC}\!\mathit{oh}}
\def\Cohstack{\mathcal{C}\!\mathit{oh}}
\def\Spacesstack{\mathcal{S}\!\mathit{paces}}
\def\Quotfunctor{\mathrm{Quot}}
\def\Hilbfunctor{\mathrm{Hilb}}
\def\Curvesstack{\mathcal{C}\!\mathit{urves}}
\def\Polarizedstack{\mathcal{P}\!\mathit{olarized}}
\def\Complexesstack{\mathcal{C}\!\mathit{omplexes}}
% \Pic is the operator that assigns to X its picard group, usage \Pic(X)
% \Picardstack_{X/B} denotes the Picard stack of X over B
% \Picardfunctor_{X/B} denotes the Picard functor of X over B
\def\Pic{\mathop{\mathrm{Pic}}\nolimits}
\def\Picardstack{\mathcal{P}\!\mathit{ic}}
\def\Picardfunctor{\mathrm{Pic}}
\def\Deformationcategory{\mathcal{D}\!\mathit{ef}}
\usepackage{cleveref}
\crefname{originalreference}{original reference}{original references}
\Crefname{originalreference}{Original reference}{Original references}
\makeatletter
\newlabel{situation-setup}{{02QL}{source}{}{sourcepointer-1}{}}
\newlabel{situation-setup@cref}{{[originalreference][0][]02QL}{[1][0][]source}{}{}{}}
\makeatother
\crefname{originalreference}{original reference}{original references}
\Crefname{originalreference}{Original reference}{Original references}
\makeatletter
\newlabel{situation-loc-chern}{{0GUJ}{source}{}{sourcepointer-2}{}}
\newlabel{situation-loc-chern@cref}{{[originalreference][0][]0GUJ}{[1][0][]source}{}{}{}}
\makeatother
\crefname{originallemma}{original lemma}{original lemmas}
\Crefname{originallemma}{Original lemma}{Original lemmas}
\makeatletter
\newlabel{lemma-independent-loc-chern-bQ}{{0FE8}{source}{}{sourcepointer-3}{}}
\newlabel{lemma-independent-loc-chern-bQ@cref}{{[originallemma][0][]0FE8}{[1][0][]source}{}{}{}}
\makeatother
\crefname{originallemma}{original lemma}{original lemmas}
\Crefname{originallemma}{Original lemma}{Original lemmas}
\makeatletter
\newlabel{lemma-silly-tensor-product}{{0FE5}{source}{}{sourcepointer-4}{}}
\newlabel{lemma-silly-tensor-product@cref}{{[originallemma][0][]0FE5}{[1][0][]source}{}{}{}}
\makeatother
\crefname{originallemma}{original lemma}{original lemmas}
\Crefname{originallemma}{Original lemma}{Original lemmas}
\makeatletter
\newlabel{lemma-homomorphism-final}{{0FE7}{source}{}{sourcepointer-5}{}}
\newlabel{lemma-homomorphism-final@cref}{{[originallemma][0][]0FE7}{[1][0][]source}{}{}{}}
\makeatother
\title{Localized Chern classes of a derived tensor product are given by the universal polynomial formula in the two localized classes}
\author{The Stacks Project Authors}
\date{stacks-project:chow.tex:local-snapshot-20261012; GNU FDL 1.2 or later}
\begin{document}
\maketitle
\noindent\textit{Editorial note.} All delivered mathematical text is the author's own, copied from the pinned Stacks Project chapter file chow.tex; no statement or proof step was rewritten, reordered or removed.\par\medskip
\noindent\textit{Editorial note.} Preamble provenance: the pinned chapter file opens with an input of the shared preamble and carries no preamble text of its own, so the project's single Stacks-wide preamble staged at sources/stacks-chow/preamble.tex is excerpted verbatim as two spans, lines 14-18 (Xy-pic setup) and lines 154-217 (the native theorem declarations and mathematical macros). Nothing was added, removed or redefined; the wrapper class and the amsmath/amssymb/amsthm/graphicx base packages are supplied by the assembler.\par\medskip
\noindent\textit{Editorial note.} The counted statement and proof contain no cross reference, no citation and no included file, so the delivered body is self contained.\par\medskip
\noindent\textit{Editorial note.} Canonical location of the counted proposition: Stacks tag 0FBF (sources/stacks-chow/pinned-original-source.tex line 11525).\par\medskip
\noindent\textit{Allowed standard prerequisites (editorial).}\begin{itemize}
\item Identification of the localized classes with those computed from a blowing up. Same chapter, Lemma labelled lemma-independent-loc-chern-bQ, cited at proof lines 11586-11592 and 11604-11618. Allows both F\_i to be handled with one common W, b and pair of complexes Q\_i.
\item The polynomial expansion of c'\_p and P'\_p for a tensor product over a common closed subscheme. Same chapter, Lemma labelled lemma-silly-tensor-product, cited at proof lines 11634-11636. Expands the localized classes of Q\_1 otimes\textasciicircum{}L Q\_2 restricted to E\_1 cap E\_2.
\item Agreement of polynomials in restricted and unrestricted localized classes. Same chapter, Lemma labelled lemma-homomorphism-final, cited at proof lines 11636-11638. Converts the identity on E\_1 cap E\_2 into the identity on Z\_1 cap Z\_2.
\item Assumption (3) of the localized Chern situation for perfect objects. Same chapter, Situation labelled situation-loc-chern, invoked in the statement. Guarantees that all the localized classes occurring in the formula are defined.
\end{itemize}\medskip
\ifdefined\hypertarget\hypertarget{sourcepointer-1}{}\fi
\noindent\textit{Source reference.} Original reference 02QL: sources/stacks-chow/pinned-original-source.tex, source lines 1372--1372. Same-chapter pointer into the author environment situation at source line 1372; invoked at proof line None of the counted proof. It is retained only as a navigation pointer; every load-bearing step is inside the delivered span.\par\medskip
\ifdefined\hypertarget\hypertarget{sourcepointer-2}{}\fi
\noindent\textit{Source reference.} Original reference 0GUJ: sources/stacks-chow/pinned-original-source.tex, source lines 10678--10678. Same-chapter pointer into the author environment situation at source line 10678; invoked at proof line None of the counted proof. It is retained only as a navigation pointer; every load-bearing step is inside the delivered span.\par\medskip
\ifdefined\hypertarget\hypertarget{sourcepointer-3}{}\fi
\noindent\textit{Source reference.} Original lemma 0FE8: sources/stacks-chow/pinned-original-source.tex, source lines 11165--11165. Same-chapter pointer into the author environment lemma at source line 11165; invoked at proof line 11571 of the counted proof. It is retained only as a navigation pointer; every load-bearing step is inside the delivered span.\par\medskip
\ifdefined\hypertarget\hypertarget{sourcepointer-4}{}\fi
\noindent\textit{Source reference.} Original lemma 0FE5: sources/stacks-chow/pinned-original-source.tex, source lines 9958--9958. Same-chapter pointer into the author environment lemma at source line 9958; invoked at proof line 11630 of the counted proof. It is retained only as a navigation pointer; every load-bearing step is inside the delivered span.\par\medskip
\ifdefined\hypertarget\hypertarget{sourcepointer-5}{}\fi
\noindent\textit{Source reference.} Original lemma 0FE7: sources/stacks-chow/pinned-original-source.tex, source lines 11026--11026. Same-chapter pointer into the author environment lemma at source line 11026; invoked at proof line 11634 of the counted proof. It is retained only as a navigation pointer; every load-bearing step is inside the delivered span.\par\medskip
\begin{lemma}
\label{lemma-loc-chern-tensor-product}
In Situation \ref{situation-setup} let $X$ be locally of finite type over $S$.
Let $Z_i \subset X$, $i = 1, 2$ be closed subschemes. Let $F_i$, $i = 1, 2$
be perfect objects of $D(\mathcal{O}_X)$. Assume for $i = 1, 2$ that
$F_i|_{X \setminus Z_i}$ is zero\footnote{Presumably there
is a variant of this lemma where we only assume $F_i|_{X \setminus Z_i}$
is isomorphic to a finite locally free $\mathcal{O}_{X \setminus Z_i}$-module
of rank $< p_i$.} and that $F_i$ on $X$ satisfies assumption
(3) of Situation \ref{situation-loc-chern}. Denote
$r_i = P_0(Z_i \to X, F_i) \in A^0(Z_i \to X)$.
Then we have
$$
c_1(Z_1 \cap Z_2 \to X, F_1 \otimes_{\mathcal{O}_X}^\mathbf{L} F_2) =
r_1 c_1(Z_2 \to X, F_2) + r_2 c_1(Z_1 \to X, F_1)
$$
in $A^1(Z_1 \cap Z_2 \to X)$ and
\begin{align*}
c_2(Z_1 \cap Z_2 \to X, F_1 \otimes_{\mathcal{O}_X}^\mathbf{L} F_2)
& =
r_1 c_2(Z_2 \to X, F_2) +
r_2 c_2(Z_1 \to X, F_1) + \\
& {r_1 \choose 2} c_1(Z_2 \to X, F_2)^2 + \\
& (r_1r_2 - 1) c_1(Z_2 \to X, F_2)c_1(Z_1 \to X, F_1) + \\
& {r_2 \choose 2} c_1(Z_1 \to X, F_1)^2
\end{align*}
in $A^2(Z_1 \cap Z_2 \to X)$ and so on for higher Chern classes.
Similarly, we have
$$
ch(Z_1 \cap Z_2 \to X, F_1 \otimes_{\mathcal{O}_X}^\mathbf{L} F_2) =
ch(Z_1 \to X, F_1) ch(Z_2 \to X, F_2)
$$
in $\prod_{p \geq 0} A^p(Z_1 \cap Z_2 \to X) \otimes \mathbf{Q}$.
More precisely, we have
$$
P_p(Z_1 \cap Z_2 \to X, F_1 \otimes_{\mathcal{O}_X}^\mathbf{L} F_2) =
\sum\nolimits_{p_1 + p_2 = p}
{p \choose p_1} P_{p_1}(Z_1 \to X, F_1) P_{p_2}(Z_2 \to X, F_2)
$$
in $A^p(Z_1 \cap Z_2 \to X)$.
\end{lemma}

\begin{proof}
Choose proper morphisms $b_i : W_i \to \mathbf{P}^1_X$ and
$Q_i \in D(\mathcal{O}_{W_i})$ as well as closed subschemes
$T_i \subset W_{i, \infty}$ as in the construction of
the localized Chern classes for $F_i$ or more generally as in
Lemma \ref{lemma-independent-loc-chern-bQ}. Choose a commutative
diagram
$$
\xymatrix{
W \ar[d]^{g_1} \ar[rd]^b \ar[r]_{g_2} & W_2 \ar[d]^{b_2} \\
W_1 \ar[r]^{b_1} & \mathbf{P}^1_X
}
$$
where all morphisms are proper and isomorphisms over
$\mathbf{A}^1_X$. For example, we can take $W$ to be the closure
of the graph of the isomorphism between
$b_1^{-1}(\mathbf{A}^1_X)$ and $b_2^{-1}(\mathbf{A}^1_X)$.
By Lemma \ref{lemma-independent-loc-chern-bQ} we may work with
$W$, $b = b_i \circ g_i$, $Lg_i^*Q_i$, and
$g_i^{-1}(T_i)$ to construct the localized Chern classes
$c_p(Z_i \to X, F_i)$. Thus we reduce to the situation described
in the next paragraph.

\medskip\noindent
Assume we have
\begin{enumerate}
\item a proper morphism $b : W \to \mathbf{P}^1_X$ which is an isomorphism
over $\mathbf{A}^1_X$,
\item $E_i \subset W_\infty$ is the inverse image of $Z_i$,
\item perfect objects $Q_i \in D(\mathcal{O}_W)$ whose Chern classes
are defined, such that
\begin{enumerate}
\item the restriction of $Q_i$ to $b^{-1}(\mathbf{A}^1_X)$ is
the pullback of $F_i$, and
\item there exists a closed subscheme $T_i \subset W_\infty$ containing
all points of $W_\infty$ lying over $X \setminus Z_i$ such that
$Q_i|_{T_i}$ is zero.
\end{enumerate}
\end{enumerate}
By Lemma \ref{lemma-independent-loc-chern-bQ} we have
$$
c_p(Z_i \to X, F_i) = c'_p(Q_i) =
(E_i \to Z_i)_* \circ c'_p(Q_i|_{E_i}) \circ C
$$
and
$$
P_p(Z_i \to X, F_i) = P'_p(Q_i) =
(E_i \to Z_i)_* \circ P'_p(Q_i|_{E_i}) \circ C
$$
for $i = 1, 2$. Next, we observe that
$Q = Q_1 \otimes_{\mathcal{O}_W}^\mathbf{L} Q_2$
satisfies (3)(a) and (3)(b) for $F_1 \otimes_{\mathcal{O}_X}^\mathbf{L} F_2$
and $T_1 \cup T_2$. Hence we see that
$$
c_p(Z_1 \cap Z_2 \to X, F_1 \otimes_{\mathcal{O}_X}^\mathbf{L} F_2) =
(E_1 \cap E_2 \to Z_1 \cap Z_2)_* \circ
c'_p(Q|_{E_1 \cap E_2}) \circ C
$$
and
$$
P_p(Z_1 \cap Z_2 \to X, F_1 \otimes_{\mathcal{O}_X}^\mathbf{L} F_2) =
(E_1 \cap E_2 \to Z_1 \cap Z_2)_* \circ
P'_p(Q|_{E_1 \cap E_2}) \circ C
$$
by the same lemma. By Lemma \ref{lemma-silly-tensor-product}
the classes $c'_p(Q|_{E_1 \cap E_2})$ and $P'_p(Q|_{E_1 \cap E_2})$
can be expanded in the correct manner in terms of the classes
$c'_p(Q_i|_{E_i})$ and $P'_p(Q_i|_{E_i})$. Then finally
Lemma \ref{lemma-homomorphism-final}
tells us that polynomials in $c'_p(Q_i|_{E_i})$ and $P'_p(Q_i|_{E_i})$
agree with the corresponding polynomials in
$c'_p(Q_i)$ and $P'_p(Q_i)$ as desired.
\end{proof}

\end{document}
